\clearpage
\thispagestyle{empty}
\SourceLeaf{417}\pdfbookmark[0]{Störungen der Pallas}{pallas}
\vspace*{1.6in}
\begin{center}{\Large STÖRUNGEN DER PALLAS.\par}\end{center}
\clearpage
\fancyhead[L]{\small STÖRUNGEN DER PALLAS}
\SourceLeaf{419}\pdfbookmark[1]{I. Briefwechsel}{pallas-briefwechsel}
\begin{center}
{\large [I. BRIEFWECHSEL.]\par}
\bigskip
\pdfbookmark[2]{Gauss an Olbers, 3. Januar 1806}{pallas-brief-18060103}
\textsc{Gauss} an \textsc{Olbers}, Braunschweig, 3. Januar 1806.
\end{center}
\noindent\hspace*{\parindent}.\quad.\quad.\quad.\quad.\quad Einige zufällige Umstände veranlassten mich vor einiger Zeit, verschiedene Untersuchungen zur Interpolationstheorie\HistoricalNoteEight{Vgl. Band III, S. 265.]} auszuarbeiten. Ich glaubte, das auf einem oder ein paar Bogen unterbringen zu können. Nun sind es aber 4 geworden, die gedruckt wohl 8 füllen könnten. Gerade die zweite, größere Hälfte brauche ich unbedingt für meine Methode zur Berechnung der Störungen und kann sie dabei sehr vorteilhaft nutzen. Daher würde ich mir wünschen, die Abhandlung vorher irgendwo drucken lassen zu können.\quad.\quad.\quad.\quad.\quad.
\par\bigskip
\begin{center}
\pdfbookmark[2]{Gauss an Olbers, 6. August 1810}{pallas-brief-18100806}
\textsc{Gauss} an \textsc{Olbers}, Göttingen, 6. August 1810.
\end{center}
\noindent\hspace*{\parindent}.\quad.\quad.\quad.\quad.\quad Ich habe angefangen, einen Aufsatz\HistoricalNoteDouble{Vgl. Band VI, S. 1.]} für die Sozietät über meine Rechnungen zur Pallas zu schreiben. Außer den numerischen Ergebnissen werden Sie darin vielleicht auch mein Verfahren, aus 4 Oppositionen eine Planetenbahn zu bestimmen (genauer gesagt die elliptischen Elemente aus 4 Längen in der Bahn), sowie die Vereinfachungen, die ich bei der Methode der kleinsten Quadrate verwende, nicht ganz uninteressant finden. Läge der Termin für die Pariser Preisfrage noch 4 Monate weiter in der Zukunft, würde ich mich jetzt vielleicht noch darum bewerben. So aber ist die Zeit zu kurz, und seit 10 Monaten fehlten mir zu jeder \SourceLeaf{420}etwas umfangreicheren zusammenhängenden Arbeit völlig die Lust und der Mut. Trotzdem denke ich nun, alle Teile der Methode, mit der sich die Pallasstörungen meiner Ansicht nach zweckmäßig berechnen lassen, in etwa 4 oder 5 einzelnen Abhandlungen für die Sozietät auszuarbeiten.
\par\bigskip
\begin{center}
\pdfbookmark[2]{Gauss an Olbers, 24. Oktober 1810}{pallas-brief-18101024}
\textsc{Gauss} an \textsc{Olbers}, 24. Oktober 1810.
\end{center}
\noindent\hspace*{\parindent}.\quad.\quad.\quad.\quad.\quad Demnächst werde ich bei der Sozietät die Arbeit\HistoricalNoteEight{Vgl. Band VI, S. 1.]} über die Pallas vorlesen, die ich Ihnen, glaube ich, schon einmal erwähnt habe. Sie werden dann einen ausführlichen Bericht darüber in unseren Anzeigen finden. Einiges Theoretische ist ebenfalls darin enthalten. Zum einen entwickle ich meine schon seit mehreren Jahren verwendete Methode, aus 4 Oppositionen eine elliptische Bahn zu bestimmen, gewissermaßen als Ergänzung zum 2. Abschnitt des 2. Buchs meiner Theoria. Zum anderen erkläre ich einen praktischen Kunstgriff, der bei der Methode der kleinsten Quadrate das sonst so beschwerliche Eliminieren erheblich verkürzt.\quad.\quad.\quad.\quad.\quad.\quad Ich werde daher vielleicht bald anfangen, die Störungen eines der neuen Planeten vollständiger zu berechnen. Wüsste ich, dass das Institut die Frist für die Preisfrage noch einmal verlängert, wäre ich nicht abgeneigt, dafür die {\AstroSymbols ⚴} zu wählen. Sonst gebührt der {\AstroSymbols ⚳} der Vortritt. Da bald die 8. {\AstroSymbols ☍} beobachtet wird, ist hier auch ein befriedigenderes Ergebnis zu erwarten. Bei der {\AstroSymbols ⚴} würde ich wohl zunächst die Elemente als veränderlich ansehen und ihre Störungen während 7 Jahren durch Quadraturen bestimmen. Dafür habe ich Formeln entworfen, die mir etwas bequemer erscheinen als die von \textsc{Laplace}. Gelingt es dann, wie zu erwarten ist, die 6 bisher beobachteten Oppositionen, die sich überhaupt nicht mehr in eine einzige Ellipse fügen wollen, gut miteinander zu vereinbaren, würde ich nach meiner schon vor 4 oder 5 Jahren entworfenen Methode die Störungen in der sonst üblichen Form berechnen, als periodische Störungen der Länge, der Breite und des Radiusvektors oder, noch besser, seines Logarithmus.\quad.\quad.\quad.\quad.\quad.
\par
